\section{Related Work}

\paragraph{LLM trading policies.}
A central problem in LLM-based trading is how to translate heterogeneous financial evidence into portfolio decisions. Memory-based approaches organise information across horizons, as in FinMem and TradingGPT, while FinCon and TradingAgents distribute analysis across specialised agents~\cite{yu2024finmem,li2023tradinggpt,yu2024fincon,xiao2024tradingagents}. FinAgent broadens the evidence available to the policy through tools~\cite{zhang2024finagent}. FLAG-Trader adapts the policy itself through reinforcement learning, extending the use of trading rewards to train financial decision policies~\cite{liu2022finrlmeta,wang2025flagtrader}. These approaches address how trading policies gather evidence and use it to select actions. When evidence processing and position selection are evaluated jointly, however, changes in portfolio performance may reflect either the information extracted from text or the exposures induced by the policy. We therefore identify an attribution gap in measuring the value of the text signal under a common exposure rule.

\paragraph{News signals and return predictability.}
Evaluating a news signal requires distinguishing information about future returns from information about price movements that have already occurred. Studies use LLM-derived sentiment scores and text embeddings to construct signals for return prediction and portfolio sorting~\cite{lopezlira2023chatgpt,kirtac2024sentiment,chen2022expected}. Work on look-ahead bias and memorisation examines a separate threat to inference, namely that a model may have encountered information about the evaluation period during pretraining~\cite{glasserman2023lookahead,lopezlira2025memorization,gao2025lookahead}. Even if pretraining overlap can be ruled out, date-only news leaves the ordering of publication and intraday price changes uncertain. For such data, we regard the incremental value of event scores as unresolved unless forward-return tests distinguish the signal from contemporaneous price movements and from proxies for news volume.

\paragraph{Benchmarking and robustness.}
The reliability of LLM trading policies is difficult to assess when results depend on the model backbone, market conditions, or the information supplied to the agent. InvestorBench facilitates comparisons under common tasks across assets and backbones, while FINSABER examines performance over longer horizons and across market conditions~\cite{li2024investorbench,li2025finsaber}. AutoRedTrader tests sensitivity to misleading financial news~\cite{liu2026autoredtrader}, while TradeTrap stress-tests components of the trading pipeline~\cite{tradetrap2025}. These evaluations address different sources of variation, but a change in strategy performance after a backbone substitution may reflect changes in both text interpretation and exposure. The remaining question is how much backbone sensitivity persists when the exposure rule is held fixed.

\paragraph{Volatility-managed portfolios.}
Changes in realised volatility can leave portfolios with substantially different levels of risk over time. Volatility management addresses this problem by reducing exposure when estimated risk rises. Prior work reports improvements in risk-adjusted performance or reductions in downside risk for some portfolios, although gains can weaken under implementable out-of-sample rules~\cite{moreira2017volmanaged,harvey2018impact,cederburg2020volmanaged}. Guidance on backtest design also highlights the risk of attributing gains to a rule selected using the evaluation data~\cite{lopezdeprado2018advances}. For a news-based strategy, we therefore ask whether the signal improves risk-adjusted performance relative to a price-only allocation subject to the same volatility rule.
